\documentclass[12pt]{article}
\usepackage[utf8]{inputenc}
\usepackage[T1]{fontenc}
\usepackage{hyperref}
\title{Uru Language Resources for K-12 Curriculum Development}
\author{Compiled by Claude Sonnet 5 (claude-sonnet-5) for Andr\'as Kornai}
\date{}
\begin{document}
\maketitle
\section{Introduction}
Uru (ISO 639-3 ure; Glottocode uruu1244) is a dormant language of Bolivia of the Uru-Chipaya family. Search results say the last fluent speaker of the Irohito variant, Julia Vila, died in 2004, and that communities around Lake Poop\'o are seeking to revive the language or its closest relative, Chipaya. Material specific to Uru is scarce.
\section{Primary Educational Resources}
An Associated Press report describes classes in which Chipaya teachers invited by the Uru communities teach numbers, songs and greetings, with radio and video used during the pandemic. No published textbook was found.
\begin{itemize}
\item AP-derived news report on the Uru and the loss of Lake Poop\'o: \url{https://www.therepublic.com/?p=450325}
\end{itemize}
\section{Linguistic and Reference Materials}
\begin{itemize}
\item DoBeS Uru-Chipaya project, people page (Irohito community and revitalization plans): \url{https://dobes.mpi.nl/projects/uru-chipaya/people/}
\item Article on Chipaya reversing predictions of extinction, Indiana journal: \url{https://journals.iai.spk-berlin.de/index.php/indiana/article/view/2486/2090}
\item CLACSO repository record: \url{https://biblioteca-repositorio.clacso.edu.ar/xmlui/handle/CLACSO/151307}
\end{itemize}
\section{Religious and Community Resources}
\begin{itemize}
\item Public Anthropologist blog on the Uru and Bolivia's 2019 intangible heritage recognition: \url{https://publicanthropologist.cmi.no/?p=1155}
\end{itemize}
\section{Digital and Audio Resources}
None found beyond the DoBeS project page, whose archive may hold recordings (not verified).
\section{Curriculum Development Context}
Revitalization would rely on Chipaya as the nearest living relative and on archival documentation. The communities' own request for classes is the strongest basis; any curriculum should be community-led.
\end{document}